% ============================================================
%  CSC343H5 Sprint 4: Normalisation
%  Student Submission Template
% ============================================================
\documentclass[11pt, letterpaper]{article}

% -- Packages ------------------------------------------------
\usepackage[margin=1in, top=0.9in, bottom=0.85in]{geometry}
\usepackage[T1]{fontenc}
\usepackage[expansion=false]{microtype}
\usepackage{xcolor}
\usepackage{enumitem}
\usepackage{parskip}
\usepackage{titlesec}
\usepackage{fancyhdr}
\usepackage{hyperref}
\usepackage{amsmath}
\usepackage{amssymb}

% -- Colours (text only, no fills) ---------------------------
\definecolor{dark}{HTML}{1A2A3A}
\definecolor{mid}{HTML}{2D5F8A}
\definecolor{muted}{HTML}{6B7A8D}

\hypersetup{colorlinks=true, linkcolor=mid, urlcolor=mid,
  pdftitle={CSC343 Sprint 4 Submission}}

% -- Section styling -----------------------------------------
\titleformat{\section}
  {\large\bfseries\color{mid}}
  {\thesection.}{0.5em}{}[\vspace{-4pt}\textcolor{mid}{\rule{\linewidth}{0.4pt}}]
\titleformat{\subsection}
  {\normalsize\bfseries\color{dark}}
  {\thesubsection}{0.5em}{}
\titlespacing{\section}{0pt}{14pt}{6pt}
\titlespacing{\subsection}{0pt}{8pt}{3pt}

% -- Header / footer -----------------------------------------
\pagestyle{fancy}
\fancyhf{}
\renewcommand{\headrulewidth}{0.4pt}
\fancyhead[L]{\small\color{muted}CSC343  Sprint 4 Submission}
\fancyhead[R]{\small\color{muted}Due: Jul 25, 5\,PM ET}
\fancyfoot[C]{\small\color{muted}\thepage}

% -- Helpers -------------------------------------------------
\newcommand{\attr}[1]{\texttt{#1}}
\newcommand{\fd}{\ensuremath{\rightarrow}}
\newcommand{\clo}[1]{\ensuremath{\{#1\}^{+}}}
\newcommand{\pk}[1]{\underline{\texttt{#1}}}          % underline a PK attribute
\newcommand{\todo}[1]{{\color{muted}\textit{[#1]}}}    % placeholder prompt

\setlist[itemize]{leftmargin=1.4em, itemsep=2pt, topsep=3pt}
\setlist[enumerate]{leftmargin=1.4em, itemsep=2pt, topsep=3pt}
\setlength{\emergencystretch}{3em}

% ============================================================
\begin{document}

% -- Title block ---------------------------------------------
\begin{center}
  {\small\color{muted}CSC343H5 Introduction to Databases}\\[4pt]
  {\LARGE\bfseries\color{dark}Sprint 4: Normalisation}\\[4pt]
  {\small\color{muted}Student Submission Template}
\end{center}
\vspace{6pt}

% -- Student information --------------------------------------
\begin{tabular}{@{}ll@{}}
  \textbf{Full name:}      & \todo{your full name} \\
  \textbf{Student number:} & \todo{your student number} \\
  \textbf{UTORid:}         & \todo{your UTORid} \\
  \textbf{Course:}         & CSC343H5 Introduction to Databases \\
  \textbf{Sprint:}         & 4: Normalisation \\
  \textbf{Submission:}     & First submission \\
\end{tabular}

\vspace{8pt}

\textbf{How to use this template.} Write each answer beneath its prompt, replacing the grey \todo{placeholder} text. Use \verb|\fd| for the arrow (\attr{A} \fd{} \attr{B}), \verb|\clo{...}| for a closure ($\clo{\attr{A}}$), and \verb|\pk{...}| to underline a primary-key attribute. Show every step; closures and decomposition steps are where the marks are.

\medskip
\textbf{Reference: \attr{research\_flat}.} Key
$(\attr{grant\_code},\ \attr{project\_id},\ \attr{researcher\_id})$. Attributes:
\attr{grant\_code}, \attr{grant\_title}, \attr{funding\_body}, \attr{funding\_sector},
\attr{project\_id}, \attr{project\_title}, \attr{budget}, \attr{researcher\_id},
\attr{researcher\_name}, \attr{institution}, \attr{role\_on\_project}. Given FDs:

\begin{itemize}[itemsep=1pt, topsep=2pt]
  \item FD1: \attr{grant\_code} \fd{} \attr{grant\_title}, \attr{funding\_body}
  \item FD2: \attr{project\_id} \fd{} \attr{project\_title}, \attr{budget}
  \item FD3: \attr{researcher\_id} \fd{} \attr{researcher\_name}, \attr{institution}
  \item FD4: \attr{funding\_body} \fd{} \attr{funding\_sector}
  \item FD5: \attr{grant\_code}, \attr{project\_id}, \attr{researcher\_id} \fd{} \attr{role\_on\_project}
\end{itemize}

% ============================================================
\section{Task 1: Closure \& Candidate Key}

\subsection*{(a) Closure of the composite key}
Let \(K = \{\attr{grant\_code}, \attr{project\_id}, \attr{researcher\_id}\}\). Show each step; the final set must contain every attribute.
\[
\begin{aligned}
K^{+} &= K && \text{(start)}\\
      &= \todo{add attributes} && \text{(by \todo{FD?})}\\
      &\ \ \ \vdots && \\
      &= \todo{all attributes} && \text{(done)}
\end{aligned}
\]

\subsection*{(b) Candidate key}
\todo{State whether the composite key is the only candidate key (or list any others), and justify in one sentence, e.g.\ which attributes never appear on a right-hand side and so must be in every key.}

\subsection*{(c) Prime and non-prime attributes}
\textbf{Prime:} \todo{...} \qquad \textbf{Non-prime:} \todo{...}

% ============================================================
\section{Task 2: Minimal Cover}

\subsection*{Step 1: Single-attribute right-hand sides}
\todo{Rewrite every FD so each has one attribute on the right.}

\subsection*{Step 2: Remove redundant left-hand-side attributes}
Only FDs with more than one attribute on the left need checking. For each, show the closure test.
\[
\clo{\todo{subset of the LHS}} = \{\,\todo{...}\,\}
\quad\Rightarrow\quad \todo{attribute is / is not redundant}
\]
\todo{Repeat as needed.}

\subsection*{Step 3: Remove redundant FDs}
For each FD $X \fd Y$, compute $\clo{X}$ \emph{without} that FD; if $Y$ is still reached, the FD is redundant.
\[
\clo{\todo{X}}\ \text{(without \todo{FD?})} = \{\,\todo{...}\,\}
\quad\Rightarrow\quad \todo{keep / drop}
\]
\todo{Repeat as needed.}

\subsection*{Resulting minimal cover}
\[ \text{Minimal cover} = \{\ \todo{list your final FDs}\ \} \]

% ============================================================
\section{Task 3: BCNF Decomposition}

For each step, name the violating FD, give the two resulting relations with their keys, and show the lossless check. Add as many steps as you need.

\subsection*{Decomposition steps}
\textbf{Step 1.} Decompose \attr{\todo{relation}} on \ \todo{$X$} \fd{} \todo{$Y$}.
\begin{itemize}
  \item \textbf{Violation:} \todo{$X$} is not a superkey of \attr{\todo{relation}} because \todo{...}.
  \item \textbf{Resulting relations:}
    \attr{\todo{R1(attrs)}} with key \pk{\todo{...}}; \quad
    \attr{\todo{R2(attrs)}} with key \pk{\todo{...}}.
  \item \textbf{Lossless check:} shared attributes $=$ \todo{...}; these form a key of \todo{which side}, so the split is lossless.
\end{itemize}

\textbf{Step 2.} \todo{...}

\textbf{Step 3.} \todo{...}

\subsection*{Final BCNF schema}
List every relation in schema notation with primary keys underlined, for example:
\begin{itemize}
  \item \attr{grant(}\pk{grant\_code}\attr{, grant\_title, funding\_body)}
  \item \todo{...}
\end{itemize}

\subsection*{BCNF verification}
\todo{For each final relation, state that every non-trivial FD in it has a superkey left-hand side.}

% ============================================================
\section{Task 4: 3NF Check \& Justification}

\subsection*{(a) Dependency preservation}
\todo{For each FD in your minimal cover, say which single relation can check it. State whether the BCNF schema is dependency-preserving. If not, give the 3NF synthesis result and say which schema you would use and why.}

\subsection*{(b) Anomaly example}
\todo{Pick one anomaly type (update, insertion, or deletion) and give a concrete example using the sample rows: which rows are involved and what goes wrong.}

% ============================================================
\section{AI Usage Declaration}

\todo{Declare any AI tool use per the course policy. If none, write: ``I did not use any AI tools in completing this sprint.''}

\begin{itemize}
  \item \textbf{Tool:} \todo{e.g.\ ChatGPT} \\
        \textbf{What I asked:} \todo{...} \\
        \textbf{How I verified / applied it:} \todo{...}
\end{itemize}

% ============================================================
\section{Self-Check (optional)}

\begin{itemize}[itemsep=1pt]
  \item[$\square$] Closure reaches all 11 attributes.
  \item[$\square$] Every minimal-cover step shows a closure test.
  \item[$\square$] Each BCNF step has a violation, two relations with keys, and a lossless check.
  \item[$\square$] Every final relation is verified BCNF.
  \item[$\square$] Name, student number, and UTORid filled in above.
\end{itemize}

\end{document}